%%%%%%%%%%%%%%%%%%%%%%%%%%%%%%%%%%%%%%%%%%%%%%%%%%%%%%%%%%%%
% 10.3 NONLINEAR PROGRAMMING
% The Composition of Advanced Mathematics
%%%%%%%%%%%%%%%%%%%%%%%%%%%%%%%%%%%%%%%%%%%%%%%%%%%%%%%%%%%%

\section{Nonlinear Programming}
\label{sec:nonlinear-programming}

\prerequisite{
Optimization fundamentals, gradients, Hessian matrices, multivariable
calculus, linear algebra, constrained optimization, Lagrange multipliers,
Karush--Kuhn--Tucker conditions, and basic numerical optimization.
}

\learningobjectives{
By the end of this section, the reader should be able to:
\begin{itemize}
    \item define a nonlinear programming problem;
    \item distinguish linear and nonlinear optimization models;
    \item identify smooth and nonsmooth nonlinear programs;
    \item formulate general constrained nonlinear optimization problems;
    \item distinguish local and global nonlinear optimization;
    \item understand feasible directions and tangent spaces;
    \item recognize active constraints;
    \item formulate the Lagrangian;
    \item derive first-order equality-constrained optimality conditions;
    \item understand constraint qualifications;
    \item state the KKT conditions for nonlinear programs;
    \item interpret complementary slackness;
    \item identify LICQ conceptually;
    \item identify MFCQ conceptually;
    \item understand second-order necessary conditions;
    \item understand second-order sufficient conditions;
    \item recognize the Hessian of the Lagrangian;
    \item understand reduced Hessian ideas conceptually;
    \item recognize saddle-point structure in KKT systems;
    \item solve simple nonlinear equality-constrained problems;
    \item solve simple nonlinear inequality-constrained problems;
    \item understand penalty methods;
    \item distinguish quadratic and exact penalties conceptually;
    \item understand barrier methods;
    \item understand logarithmic barriers;
    \item understand augmented Lagrangian methods;
    \item understand sequential quadratic programming;
    \item understand nonlinear interior-point methods conceptually;
    \item recognize trust-region methods for constrained optimization;
    \item understand line-search globalization;
    \item understand merit functions;
    \item recognize feasibility restoration;
    \item understand Newton systems for KKT equations;
    \item recognize sparse nonlinear programming structure;
    \item understand scaling and conditioning;
    \item distinguish local solvers from global nonlinear optimization;
    \item recognize nonconvexity and multiple local optima;
    \item understand multistart strategies conceptually;
    \item connect nonlinear programming with engineering design,
          estimation, economics, machine learning, and optimal control.
\end{itemize}
}

%%%%%%%%%%%%%%%%%%%%%%%%%%%%%%%%%%%%%%%%%%%%%%%%%%%%%%%%%%%%
\subsection{What Is Nonlinear Programming?}
\label{subsec:what-is-nonlinear-programming}
%%%%%%%%%%%%%%%%%%%%%%%%%%%%%%%%%%%%%%%%%%%%%%%%%%%%%%%%%%%%

A nonlinear programming problem, abbreviated NLP, is an optimization
problem in which the objective function, one or more constraints, or both
are nonlinear.

A general form is

\[
\boxed{
\begin{aligned}
\min_{\mathbf{x}\in\mathbb{R}^n}
\quad&
f(\mathbf{x})
\\
\text{subject to}
\quad&
g_i(\mathbf{x})\leq0,
\qquad
i=1,\ldots,m,
\\
&
h_j(\mathbf{x})=0,
\qquad
j=1,\ldots,p.
\end{aligned}
}
\]

The functions \(f\), \(g_i\), and \(h_j\) may be nonlinear.

%%%%%%%%%%%%%%%%%%%%%%%%%%%%%%%%%%%%%%%%%%%%%%%%%%%%%%%%%%%%
\subsection{Linear Versus Nonlinear Programming}
\label{subsec:linear-versus-nonlinear-programming}
%%%%%%%%%%%%%%%%%%%%%%%%%%%%%%%%%%%%%%%%%%%%%%%%%%%%%%%%%%%%

A linear program requires:

\[
\boxed{
\text{linear objective}
}
\]

and

\[
\boxed{
\text{linear constraints}.
}
\]

A nonlinear program may contain terms such as

\[
x_1^2,
\]

\[
x_1x_2,
\]

\[
e^{x_1},
\]

\[
\log x_2,
\]

or

\[
\sin x_1.
\]

For example,

\[
\boxed{
\min_{x,y}
x^2+y^2
}
\]

subject to

\[
\boxed{
xy\geq1
}
\]

is nonlinear because the constraint contains the product \(xy\).

%%%%%%%%%%%%%%%%%%%%%%%%%%%%%%%%%%%%%%%%%%%%%%%%%%%%%%%%%%%%
\subsection{Smooth Nonlinear Programs}
\label{subsec:smooth-nonlinear-programs}
%%%%%%%%%%%%%%%%%%%%%%%%%%%%%%%%%%%%%%%%%%%%%%%%%%%%%%%%%%%%

A smooth nonlinear program has sufficiently differentiable objective and
constraint functions.

Typical algorithms rely on:

\[
\boxed{
\nabla f,
}
\]

\[
\boxed{
\nabla g_i,
}
\]

\[
\boxed{
\nabla h_j,
}
\]

and sometimes

\[
\boxed{
\nabla^2 f.
}
\]

Smoothness enables Taylor expansion and local optimality analysis.

%%%%%%%%%%%%%%%%%%%%%%%%%%%%%%%%%%%%%%%%%%%%%%%%%%%%%%%%%%%%
\subsection{Nonsmooth Nonlinear Programs}
\label{subsec:nonsmooth-nonlinear-programs}
%%%%%%%%%%%%%%%%%%%%%%%%%%%%%%%%%%%%%%%%%%%%%%%%%%%%%%%%%%%%

A nonlinear optimization problem can also be nonsmooth.

For example,

\[
\boxed{
\min_x
|x|
}
\]

has no ordinary derivative at

\[
x=0.
\]

Similarly,

\[
\boxed{
\min_{\mathbf{x}}
f(\mathbf{x})
+
\lambda
\|\mathbf{x}\|_1
}
\]

is nonsmooth at points where components vanish.

Nonsmooth optimization requires generalized derivative concepts and
specialized methods.

%%%%%%%%%%%%%%%%%%%%%%%%%%%%%%%%%%%%%%%%%%%%%%%%%%%%%%%%%%%%
\subsection{Local and Global Solutions}
\label{subsec:nlp-local-global-solutions}
%%%%%%%%%%%%%%%%%%%%%%%%%%%%%%%%%%%%%%%%%%%%%%%%%%%%%%%%%%%%

A feasible point \(\mathbf{x}^*\) is a local minimizer if there exists a
neighborhood in which

\[
\boxed{
f(\mathbf{x}^*)
\leq
f(\mathbf{x})
}
\]

for all nearby feasible \(\mathbf{x}\).

A global minimizer satisfies

\[
\boxed{
f(\mathbf{x}^*)
\leq
f(\mathbf{x})
}
\]

for every feasible point.

For nonlinear and especially nonconvex problems,

\[
\boxed{
\text{local minimum}
\not\Rightarrow
\text{global minimum}.
}
\]

%%%%%%%%%%%%%%%%%%%%%%%%%%%%%%%%%%%%%%%%%%%%%%%%%%%%%%%%%%%%
\subsection{Nonconvexity}
\label{subsec:nlp-nonconvexity}
%%%%%%%%%%%%%%%%%%%%%%%%%%%%%%%%%%%%%%%%%%%%%%%%%%%%%%%%%%%%

A nonlinear problem is nonconvex if the objective or feasible region fails
to have the required convex structure.

Nonconvexity can produce:

\[
\boxed{
\text{multiple local minima},
}
\]

\[
\boxed{
\text{local maxima},
}
\]

\[
\boxed{
\text{saddle points},
}
\]

and

\[
\boxed{
\text{disconnected feasible regions}.
}
\]

This makes global optimization difficult.

%%%%%%%%%%%%%%%%%%%%%%%%%%%%%%%%%%%%%%%%%%%%%%%%%%%%%%%%%%%%
\subsection{Example of Multiple Local Minima}
\label{subsec:nlp-multiple-local-minima}
%%%%%%%%%%%%%%%%%%%%%%%%%%%%%%%%%%%%%%%%%%%%%%%%%%%%%%%%%%%%

Consider

\[
\boxed{
f(x)
=
(x^2-1)^2.
}
\]

Then

\[
f'(x)
=
4x(x^2-1).
\]

Stationary points are

\[
x=-1,
\qquad
x=0,
\qquad
x=1.
\]

The points

\[
\boxed{
x=\pm1
}
\]

are global minima, while

\[
\boxed{
x=0
}
\]

is a local maximum.

%%%%%%%%%%%%%%%%%%%%%%%%%%%%%%%%%%%%%%%%%%%%%%%%%%%%%%%%%%%%
\subsection{Feasible Directions}
\label{subsec:nlp-feasible-directions}
%%%%%%%%%%%%%%%%%%%%%%%%%%%%%%%%%%%%%%%%%%%%%%%%%%%%%%%%%%%%

Let \(\mathbf{x}\) be feasible.

A direction \(\mathbf{d}\) is a feasible direction if sufficiently small
movement in that direction remains feasible:

\[
\boxed{
\mathbf{x}
+
\alpha\mathbf{d}
\in
\mathcal{F}
}
\]

for sufficiently small positive \(\alpha\).

At boundary points, not every direction is feasible.

%%%%%%%%%%%%%%%%%%%%%%%%%%%%%%%%%%%%%%%%%%%%%%%%%%%%%%%%%%%%
\subsection{Tangent Directions for Equality Constraints}
\label{subsec:nlp-tangent-directions}
%%%%%%%%%%%%%%%%%%%%%%%%%%%%%%%%%%%%%%%%%%%%%%%%%%%%%%%%%%%%

Suppose

\[
h_j(\mathbf{x})=0.
\]

A tangent direction \(\mathbf{d}\) at \(\mathbf{x}^*\) satisfies

\[
\boxed{
\nabla h_j(\mathbf{x}^*)^T
\mathbf{d}
=
0
}
\]

for all equality constraints.

Thus feasible first-order motion lies in the tangent space.

%%%%%%%%%%%%%%%%%%%%%%%%%%%%%%%%%%%%%%%%%%%%%%%%%%%%%%%%%%%%
\subsection{Active Inequality Constraints}
\label{subsec:nlp-active-constraints}
%%%%%%%%%%%%%%%%%%%%%%%%%%%%%%%%%%%%%%%%%%%%%%%%%%%%%%%%%%%%

For

\[
g_i(\mathbf{x})\leq0,
\]

constraint \(i\) is active at \(\mathbf{x}^*\) if

\[
\boxed{
g_i(\mathbf{x}^*)=0.
}
\]

It is inactive if

\[
\boxed{
g_i(\mathbf{x}^*)<0.
}
\]

Only active inequalities contribute to the local boundary geometry.

%%%%%%%%%%%%%%%%%%%%%%%%%%%%%%%%%%%%%%%%%%%%%%%%%%%%%%%%%%%%
\subsection{Equality-Constrained Nonlinear Optimization}
\label{subsec:nlp-equality-constrained}
%%%%%%%%%%%%%%%%%%%%%%%%%%%%%%%%%%%%%%%%%%%%%%%%%%%%%%%%%%%%

Consider

\[
\boxed{
\begin{aligned}
\min_{\mathbf{x}}
\quad&
f(\mathbf{x})
\\
\text{subject to}
\quad&
h_j(\mathbf{x})=0,
\qquad
j=1,\ldots,p.
\end{aligned}
}
\]

The Lagrangian is

\[
\boxed{
\mathcal{L}
(
\mathbf{x},
\boldsymbol{\lambda}
)
=
f(\mathbf{x})
+
\sum_{j=1}^{p}
\lambda_j
h_j(\mathbf{x}).
}
\]

%%%%%%%%%%%%%%%%%%%%%%%%%%%%%%%%%%%%%%%%%%%%%%%%%%%%%%%%%%%%
\subsection{First-Order Equality Conditions}
\label{subsec:nlp-first-order-equality}
%%%%%%%%%%%%%%%%%%%%%%%%%%%%%%%%%%%%%%%%%%%%%%%%%%%%%%%%%%%%

Under appropriate regularity conditions, a local constrained optimum
satisfies

\[
\boxed{
\nabla_{\mathbf{x}}
\mathcal{L}
(
\mathbf{x}^*,
\boldsymbol{\lambda}^*
)
=
0,
}
\]

together with

\[
\boxed{
h_j(\mathbf{x}^*)=0.
}
\]

Thus,

\[
\boxed{
\nabla f(\mathbf{x}^*)
+
\sum_{j=1}^{p}
\lambda_j^*
\nabla h_j(\mathbf{x}^*)
=
0.
}
\]

%%%%%%%%%%%%%%%%%%%%%%%%%%%%%%%%%%%%%%%%%%%%%%%%%%%%%%%%%%%%
\subsection{Worked Example: Nonlinear Equality Constraint}
\label{subsec:worked-nlp-equality}
%%%%%%%%%%%%%%%%%%%%%%%%%%%%%%%%%%%%%%%%%%%%%%%%%%%%%%%%%%%%

\begin{workedexamplebox}{Minimum on a Circle}

Minimize

\[
f(x,y)=x+y
\]

subject to

\[
x^2+y^2=1.
\]

The Lagrangian is

\[
\mathcal{L}
=
x+y
+
\lambda
(
x^2+y^2-1
).
\]

Stationarity gives

\[
1+2\lambda x=0,
\]

\[
1+2\lambda y=0.
\]

Therefore,

\[
x=y.
\]

Using the constraint,

\[
2x^2=1.
\]

Thus,

\[
x=y=\pm\frac1{\sqrt2}.
\]

Evaluate the objective:

\[
f
\left(
\frac1{\sqrt2},
\frac1{\sqrt2}
\right)
=
\sqrt2,
\]

while

\[
f
\left(
-\frac1{\sqrt2},
-\frac1{\sqrt2}
\right)
=
-\sqrt2.
\]

Therefore,

\begin{answerbox}
\[
\boxed{
(x^*,y^*)
=
\left(
-\frac1{\sqrt2},
-\frac1{\sqrt2}
\right),
}
\]

with minimum value

\[
\boxed{
f^*=-\sqrt2.
}
\]
\end{answerbox}

\end{workedexamplebox}

%%%%%%%%%%%%%%%%%%%%%%%%%%%%%%%%%%%%%%%%%%%%%%%%%%%%%%%%%%%%
\subsection{Constraint Qualifications}
\label{subsec:nlp-constraint-qualifications}
%%%%%%%%%%%%%%%%%%%%%%%%%%%%%%%%%%%%%%%%%%%%%%%%%%%%%%%%%%%%

Constraint qualifications are regularity conditions ensuring that local
constraint geometry is represented correctly by the gradients.

They are needed because multiplier conditions can fail at degenerate
points.

Important examples include:

\[
\boxed{
\text{LICQ},
}
\]

\[
\boxed{
\text{MFCQ},
}
\]

and

\[
\boxed{
\text{Slater's condition}
}
\]

for convex problems.

%%%%%%%%%%%%%%%%%%%%%%%%%%%%%%%%%%%%%%%%%%%%%%%%%%%%%%%%%%%%
\subsection{Linear Independence Constraint Qualification}
\label{subsec:licq}
%%%%%%%%%%%%%%%%%%%%%%%%%%%%%%%%%%%%%%%%%%%%%%%%%%%%%%%%%%%%

LICQ stands for Linear Independence Constraint Qualification.

At \(\mathbf{x}^*\), LICQ requires the gradients of:

\[
\boxed{
\text{all equality constraints}
}
\]

together with

\[
\boxed{
\text{all active inequality constraints}
}
\]

to be linearly independent.

Under LICQ, KKT multipliers are locally well behaved and are unique under
appropriate conditions.

%%%%%%%%%%%%%%%%%%%%%%%%%%%%%%%%%%%%%%%%%%%%%%%%%%%%%%%%%%%%
\subsection{Mangasarian--Fromovitz Constraint Qualification}
\label{subsec:mfcq}
%%%%%%%%%%%%%%%%%%%%%%%%%%%%%%%%%%%%%%%%%%%%%%%%%%%%%%%%%%%%

MFCQ is weaker than LICQ.

Conceptually, it requires:

\begin{itemize}
    \item equality gradients to be linearly independent; and
    \item existence of a direction tangent to the equality constraints
          that strictly decreases all active inequality constraints.
\end{itemize}

MFCQ ensures the feasible set has an appropriate local direction structure.

%%%%%%%%%%%%%%%%%%%%%%%%%%%%%%%%%%%%%%%%%%%%%%%%%%%%%%%%%%%%
\subsection{Failure of Constraint Qualification}
\label{subsec:nlp-cq-failure}
%%%%%%%%%%%%%%%%%%%%%%%%%%%%%%%%%%%%%%%%%%%%%%%%%%%%%%%%%%%%

Consider

\[
\boxed{
\min_x x
}
\]

subject to

\[
\boxed{
x^2=0.
}
\]

The only feasible point is

\[
x=0,
\]

so it is trivially optimal.

However,

\[
\nabla h(0)=0.
\]

The ordinary Lagrange condition would require

\[
1+\lambda(0)=0,
\]

which is impossible.

The issue is failure of constraint regularity.

%%%%%%%%%%%%%%%%%%%%%%%%%%%%%%%%%%%%%%%%%%%%%%%%%%%%%%%%%%%%
\subsection{General Nonlinear KKT Conditions}
\label{subsec:nlp-kkt}
%%%%%%%%%%%%%%%%%%%%%%%%%%%%%%%%%%%%%%%%%%%%%%%%%%%%%%%%%%%%

Consider

\[
\boxed{
\begin{aligned}
\min_{\mathbf{x}}
\quad&
f(\mathbf{x})
\\
\text{subject to}
\quad&
g_i(\mathbf{x})\leq0,
\qquad i=1,\ldots,m,
\\
&
h_j(\mathbf{x})=0,
\qquad j=1,\ldots,p.
\end{aligned}
}
\]

Define

\[
\boxed{
\mathcal{L}
=
f
+
\sum_{i=1}^{m}
\mu_i g_i
+
\sum_{j=1}^{p}
\lambda_j h_j.
}
\]

At a regular local minimizer, the KKT conditions hold.

%%%%%%%%%%%%%%%%%%%%%%%%%%%%%%%%%%%%%%%%%%%%%%%%%%%%%%%%%%%%
\subsection{KKT Stationarity}
\label{subsec:nlp-kkt-stationarity}
%%%%%%%%%%%%%%%%%%%%%%%%%%%%%%%%%%%%%%%%%%%%%%%%%%%%%%%%%%%%

Stationarity requires

\[
\boxed{
\nabla f(\mathbf{x}^*)
+
\sum_{i=1}^{m}
\mu_i^*
\nabla g_i(\mathbf{x}^*)
+
\sum_{j=1}^{p}
\lambda_j^*
\nabla h_j(\mathbf{x}^*)
=
0.
}
\]

%%%%%%%%%%%%%%%%%%%%%%%%%%%%%%%%%%%%%%%%%%%%%%%%%%%%%%%%%%%%
\subsection{KKT Primal Feasibility}
\label{subsec:nlp-kkt-primal-feasibility}
%%%%%%%%%%%%%%%%%%%%%%%%%%%%%%%%%%%%%%%%%%%%%%%%%%%%%%%%%%%%

The candidate must satisfy

\[
\boxed{
g_i(\mathbf{x}^*)\leq0,
}
\]

and

\[
\boxed{
h_j(\mathbf{x}^*)=0.
}
\]

%%%%%%%%%%%%%%%%%%%%%%%%%%%%%%%%%%%%%%%%%%%%%%%%%%%%%%%%%%%%
\subsection{KKT Dual Feasibility}
\label{subsec:nlp-kkt-dual-feasibility}
%%%%%%%%%%%%%%%%%%%%%%%%%%%%%%%%%%%%%%%%%%%%%%%%%%%%%%%%%%%%

For the convention

\[
g_i(\mathbf{x})\leq0,
\]

the multipliers satisfy

\[
\boxed{
\mu_i^*\geq0.
}
\]

Equality multipliers remain unrestricted in sign.

%%%%%%%%%%%%%%%%%%%%%%%%%%%%%%%%%%%%%%%%%%%%%%%%%%%%%%%%%%%%
\subsection{KKT Complementary Slackness}
\label{subsec:nlp-kkt-complementary-slackness}
%%%%%%%%%%%%%%%%%%%%%%%%%%%%%%%%%%%%%%%%%%%%%%%%%%%%%%%%%%%%

For every inequality,

\[
\boxed{
\mu_i^*
g_i(\mathbf{x}^*)
=
0.
}
\]

Hence,

\[
g_i(\mathbf{x}^*)<0
\Rightarrow
\mu_i^*=0.
\]

Only active constraints can have positive multipliers.

%%%%%%%%%%%%%%%%%%%%%%%%%%%%%%%%%%%%%%%%%%%%%%%%%%%%%%%%%%%%
\subsection{Worked Example: Nonlinear Inequality Constraint}
\label{subsec:worked-nlp-inequality}
%%%%%%%%%%%%%%%%%%%%%%%%%%%%%%%%%%%%%%%%%%%%%%%%%%%%%%%%%%%%

\begin{workedexamplebox}{KKT with an Active Constraint}

Minimize

\[
f(x,y)=x^2+y^2
\]

subject to

\[
x+y\geq2.
\]

Rewrite the inequality as

\[
g(x,y)
=
2-x-y
\leq0.
\]

The Lagrangian is

\[
\mathcal{L}
=
x^2+y^2
+
\mu(2-x-y).
\]

Stationarity gives

\[
2x-\mu=0,
\]

\[
2y-\mu=0.
\]

Thus,

\[
x=y.
\]

The unconstrained minimizer \((0,0)\) is infeasible, so the constraint is
active:

\[
x+y=2.
\]

Hence,

\[
x=y=1.
\]

Then

\[
\mu=2.
\]

Therefore,

\begin{answerbox}
\[
\boxed{
(x^*,y^*)=(1,1),
}
\]

with

\[
\boxed{
\mu^*=2.
}
\]
\end{answerbox}

\end{workedexamplebox}

%%%%%%%%%%%%%%%%%%%%%%%%%%%%%%%%%%%%%%%%%%%%%%%%%%%%%%%%%%%%
\subsection{Necessary Versus Sufficient KKT Conditions}
\label{subsec:kkt-necessary-sufficient}
%%%%%%%%%%%%%%%%%%%%%%%%%%%%%%%%%%%%%%%%%%%%%%%%%%%%%%%%%%%%

For general nonlinear problems, KKT conditions are usually necessary
under a suitable constraint qualification.

They are not generally sufficient.

A KKT point may correspond to:

\[
\boxed{
\text{local minimum},
}
\]

\[
\boxed{
\text{local maximum},
}
\]

or

\[
\boxed{
\text{saddle-type behavior}.
}
\]

Convexity changes this situation substantially.

%%%%%%%%%%%%%%%%%%%%%%%%%%%%%%%%%%%%%%%%%%%%%%%%%%%%%%%%%%%%
\subsection{Hessian of the Lagrangian}
\label{subsec:hessian-of-lagrangian}
%%%%%%%%%%%%%%%%%%%%%%%%%%%%%%%%%%%%%%%%%%%%%%%%%%%%%%%%%%%%

The Hessian with respect to the decision variables is

\[
\boxed{
\nabla_{\mathbf{x}\mathbf{x}}^2
\mathcal{L}
=
\nabla^2f
+
\sum_i
\mu_i
\nabla^2g_i
+
\sum_j
\lambda_j
\nabla^2h_j.
}
\]

Second-order constrained optimality depends on this matrix restricted to
appropriate feasible directions.

%%%%%%%%%%%%%%%%%%%%%%%%%%%%%%%%%%%%%%%%%%%%%%%%%%%%%%%%%%%%
\subsection{Second-Order Necessary Condition}
\label{subsec:nlp-second-order-necessary}
%%%%%%%%%%%%%%%%%%%%%%%%%%%%%%%%%%%%%%%%%%%%%%%%%%%%%%%%%%%%

At a sufficiently regular local minimizer satisfying KKT conditions, the
Hessian of the Lagrangian must be nonnegative along appropriate critical
directions.

In a pure equality-constrained problem, if \(\mathbf{d}\) satisfies

\[
\boxed{
\nabla h_j(\mathbf{x}^*)^T\mathbf{d}=0
}
\]

for every constraint, then a second-order necessary condition is

\[
\boxed{
\mathbf{d}^T
\nabla_{\mathbf{x}\mathbf{x}}^2
\mathcal{L}
(\mathbf{x}^*,\boldsymbol{\lambda}^*)
\mathbf{d}
\geq0.
}
\]

%%%%%%%%%%%%%%%%%%%%%%%%%%%%%%%%%%%%%%%%%%%%%%%%%%%%%%%%%%%%
\subsection{Second-Order Sufficient Condition}
\label{subsec:nlp-second-order-sufficient}
%%%%%%%%%%%%%%%%%%%%%%%%%%%%%%%%%%%%%%%%%%%%%%%%%%%%%%%%%%%%

If the KKT conditions hold and

\[
\boxed{
\mathbf{d}^T
\nabla_{\mathbf{x}\mathbf{x}}^2
\mathcal{L}
\mathbf{d}
>0
}
\]

for every nonzero feasible critical direction \(\mathbf{d}\), then under
appropriate regularity conditions the point is a strict local minimizer.

Thus curvature is checked only along directions permitted by the
constraints.

%%%%%%%%%%%%%%%%%%%%%%%%%%%%%%%%%%%%%%%%%%%%%%%%%%%%%%%%%%%%
\subsection{Reduced Hessian}
\label{subsec:reduced-hessian}
%%%%%%%%%%%%%%%%%%%%%%%%%%%%%%%%%%%%%%%%%%%%%%%%%%%%%%%%%%%%

Suppose the columns of \(Z\) form a basis for the tangent space:

\[
\boxed{
\nabla h(\mathbf{x}^*)^TZ=0.
}
\]

The reduced Hessian is

\[
\boxed{
Z^T
\nabla_{\mathbf{x}\mathbf{x}}^2
\mathcal{L}
Z.
}
\]

Positive definiteness of the reduced Hessian gives a convenient
second-order test for equality-constrained minima.

%%%%%%%%%%%%%%%%%%%%%%%%%%%%%%%%%%%%%%%%%%%%%%%%%%%%%%%%%%%%
\subsection{Newton System for Equality-Constrained Optimization}
\label{subsec:equality-kkt-newton-system}
%%%%%%%%%%%%%%%%%%%%%%%%%%%%%%%%%%%%%%%%%%%%%%%%%%%%%%%%%%%%

For equality constraints,

\[
h(\mathbf{x})=0,
\]

the KKT equations are

\[
\boxed{
\nabla f(\mathbf{x})
+
J_h(\mathbf{x})^T\boldsymbol{\lambda}
=
0,
}
\]

and

\[
\boxed{
h(\mathbf{x})=0.
}
\]

Newton's method applied to these equations gives a linear system of the
form

\[
\boxed{
\begin{pmatrix}
H & J_h^T\\
J_h & 0
\end{pmatrix}
\begin{pmatrix}
\mathbf{p}\\
\Delta\boldsymbol{\lambda}
\end{pmatrix}
=
-
\begin{pmatrix}
\nabla_{\mathbf{x}}\mathcal{L}\\
h
\end{pmatrix},
}
\]

where

\[
\boxed{
H
=
\nabla_{\mathbf{x}\mathbf{x}}^2\mathcal{L}.
}
\]

%%%%%%%%%%%%%%%%%%%%%%%%%%%%%%%%%%%%%%%%%%%%%%%%%%%%%%%%%%%%
\subsection{Saddle-Point Linear Systems}
\label{subsec:nlp-saddle-point-system}
%%%%%%%%%%%%%%%%%%%%%%%%%%%%%%%%%%%%%%%%%%%%%%%%%%%%%%%%%%%%

The KKT matrix

\[
\boxed{
\begin{pmatrix}
H&A^T\\
A&0
\end{pmatrix}
}
\]

is called a saddle-point matrix.

Such matrices arise throughout:

\[
\boxed{
\text{constrained optimization},
}
\]

\[
\boxed{
\text{fluid mechanics},
}
\]

\[
\boxed{
\text{finite-element methods},
}
\]

and

\[
\boxed{
\text{optimal control}.
}
\]

%%%%%%%%%%%%%%%%%%%%%%%%%%%%%%%%%%%%%%%%%%%%%%%%%%%%%%%%%%%%
\subsection{Penalty Methods}
\label{subsec:nlp-penalty-methods}
%%%%%%%%%%%%%%%%%%%%%%%%%%%%%%%%%%%%%%%%%%%%%%%%%%%%%%%%%%%%

Penalty methods transform constrained optimization into a sequence of
unconstrained or simpler problems.

For equality constraints

\[
h(\mathbf{x})=0,
\]

a quadratic penalty objective is

\[
\boxed{
P_\rho(\mathbf{x})
=
f(\mathbf{x})
+
\frac{\rho}{2}
\|
h(\mathbf{x})
\|^2.
}
\]

As \(\rho\) increases, constraint violation becomes increasingly
expensive.

%%%%%%%%%%%%%%%%%%%%%%%%%%%%%%%%%%%%%%%%%%%%%%%%%%%%%%%%%%%%
\subsection{Quadratic Penalty Method}
\label{subsec:quadratic-penalty}
%%%%%%%%%%%%%%%%%%%%%%%%%%%%%%%%%%%%%%%%%%%%%%%%%%%%%%%%%%%%

A sequence of problems is solved:

\[
\boxed{
\min_{\mathbf{x}}
P_{\rho_k}(\mathbf{x}),
}
\]

with

\[
\boxed{
\rho_k\to\infty.
}
\]

Under suitable conditions, minimizers approach feasible solutions of the
original constrained problem.

However, very large penalty parameters can create severe ill-conditioning.

%%%%%%%%%%%%%%%%%%%%%%%%%%%%%%%%%%%%%%%%%%%%%%%%%%%%%%%%%%%%
\subsection{Worked Example: Quadratic Penalty}
\label{subsec:worked-quadratic-penalty}
%%%%%%%%%%%%%%%%%%%%%%%%%%%%%%%%%%%%%%%%%%%%%%%%%%%%%%%%%%%%

\begin{workedexamplebox}{Penalizing an Equality Constraint}

Consider

\[
\min_x
(x-2)^2
\]

subject to

\[
x=1.
\]

The constrained optimum is clearly

\[
x^*=1.
\]

The quadratic penalty objective is

\[
P_\rho(x)
=
(x-2)^2
+
\frac{\rho}{2}
(x-1)^2.
\]

Differentiate:

\[
P_\rho'(x)
=
2(x-2)
+
\rho(x-1).
\]

Set equal to zero:

\[
(2+\rho)x
=
4+\rho.
\]

Thus,

\[
\boxed{
x_\rho
=
\frac{
\rho+4
}{
\rho+2
}.
}
\]

As

\[
\rho\to\infty,
\]

\[
\begin{answerbox}
\[
\boxed{
x_\rho\to1.
}
\]
\end{answerbox}

\end{workedexamplebox}

%%%%%%%%%%%%%%%%%%%%%%%%%%%%%%%%%%%%%%%%%%%%%%%%%%%%%%%%%%%%
\subsection{Exact Penalty Functions}
\label{subsec:exact-penalty-functions}
%%%%%%%%%%%%%%%%%%%%%%%%%%%%%%%%%%%%%%%%%%%%%%%%%%%%%%%%%%%%

An exact penalty can recover a constrained solution using a finite
penalty parameter.

A common form for equality constraints is

\[
\boxed{
P_\rho(\mathbf{x})
=
f(\mathbf{x})
+
\rho
\sum_j
|h_j(\mathbf{x})|.
}
\]

Unlike quadratic penalties, exact penalties are generally nonsmooth.

%%%%%%%%%%%%%%%%%%%%%%%%%%%%%%%%%%%%%%%%%%%%%%%%%%%%%%%%%%%%
\subsection{Barrier Methods}
\label{subsec:nlp-barrier-methods}
%%%%%%%%%%%%%%%%%%%%%%%%%%%%%%%%%%%%%%%%%%%%%%%%%%%%%%%%%%%%

Barrier methods enforce inequality feasibility from the interior.

For

\[
g_i(\mathbf{x})<0,
\]

a logarithmic barrier problem is

\[
\boxed{
\min_{\mathbf{x}}
f(\mathbf{x})
-
\mu
\sum_{i=1}^{m}
\log
[
-g_i(\mathbf{x})
].
}
\]

The parameter

\[
\mu>0
\]

controls the barrier strength.

%%%%%%%%%%%%%%%%%%%%%%%%%%%%%%%%%%%%%%%%%%%%%%%%%%%%%%%%%%%%
\subsection{Barrier Behavior}
\label{subsec:nlp-barrier-behavior}
%%%%%%%%%%%%%%%%%%%%%%%%%%%%%%%%%%%%%%%%%%%%%%%%%%%%%%%%%%%%

As a feasible point approaches the boundary,

\[
g_i(\mathbf{x})\to0^-,
\]

we have

\[
-\log[-g_i(\mathbf{x})]
\to\infty.
\]

Therefore the barrier prevents iterates from crossing the inequality
boundary.

As

\[
\mu\to0,
\]

the barrier problem approaches the original constrained problem.

%%%%%%%%%%%%%%%%%%%%%%%%%%%%%%%%%%%%%%%%%%%%%%%%%%%%%%%%%%%%
\subsection{Central Path in Nonlinear Optimization}
\label{subsec:nlp-central-path}
%%%%%%%%%%%%%%%%%%%%%%%%%%%%%%%%%%%%%%%%%%%%%%%%%%%%%%%%%%%%

The solutions of the barrier problems for varying \(\mu\) form a central
path under appropriate conditions.

This path satisfies perturbed KKT conditions.

As

\[
\boxed{
\mu\to0,
}
\]

the solution approaches a constrained optimum.

%%%%%%%%%%%%%%%%%%%%%%%%%%%%%%%%%%%%%%%%%%%%%%%%%%%%%%%%%%%%
\subsection{Augmented Lagrangian Method}
\label{subsec:augmented-lagrangian}
%%%%%%%%%%%%%%%%%%%%%%%%%%%%%%%%%%%%%%%%%%%%%%%%%%%%%%%%%%%%

For equality constraints,

\[
h(\mathbf{x})=0,
\]

the augmented Lagrangian is

\[
\boxed{
\mathcal{L}_\rho
(
\mathbf{x},
\boldsymbol{\lambda}
)
=
f(\mathbf{x})
+
\boldsymbol{\lambda}^T
h(\mathbf{x})
+
\frac{\rho}{2}
\|
h(\mathbf{x})
\|^2.
}
\]

It combines multiplier information with a quadratic penalty.

%%%%%%%%%%%%%%%%%%%%%%%%%%%%%%%%%%%%%%%%%%%%%%%%%%%%%%%%%%%%
\subsection{Why Augmented Lagrangians Help}
\label{subsec:why-augmented-lagrangian}
%%%%%%%%%%%%%%%%%%%%%%%%%%%%%%%%%%%%%%%%%%%%%%%%%%%%%%%%%%%%

Pure quadratic penalty methods often require

\[
\rho\to\infty.
\]

Augmented Lagrangian methods update both:

\[
\boxed{
\boldsymbol{\lambda}
}
\]

and

\[
\boxed{
\rho.
}
\]

This can achieve accurate feasibility without requiring extremely large
penalty parameters.

%%%%%%%%%%%%%%%%%%%%%%%%%%%%%%%%%%%%%%%%%%%%%%%%%%%%%%%%%%%%
\subsection{Multiplier Update}
\label{subsec:augmented-lagrangian-multiplier-update}
%%%%%%%%%%%%%%%%%%%%%%%%%%%%%%%%%%%%%%%%%%%%%%%%%%%%%%%%%%%%

A basic equality-constrained multiplier update has form

\[
\boxed{
\boldsymbol{\lambda}_{k+1}
=
\boldsymbol{\lambda}_k
+
\rho_k
h(\mathbf{x}_{k+1}).
}
\]

The multiplier accumulates information about persistent constraint
violation.

%%%%%%%%%%%%%%%%%%%%%%%%%%%%%%%%%%%%%%%%%%%%%%%%%%%%%%%%%%%%
\subsection{Sequential Quadratic Programming}
\label{subsec:sequential-quadratic-programming}
%%%%%%%%%%%%%%%%%%%%%%%%%%%%%%%%%%%%%%%%%%%%%%%%%%%%%%%%%%%%

Sequential quadratic programming, abbreviated SQP, solves a sequence of
quadratic approximations to the nonlinear problem.

At iteration \(k\), the objective is approximated quadratically and the
constraints are linearized.

This produces a quadratic programming subproblem.

%%%%%%%%%%%%%%%%%%%%%%%%%%%%%%%%%%%%%%%%%%%%%%%%%%%%%%%%%%%%
\subsection{SQP Subproblem}
\label{subsec:sqp-subproblem}
%%%%%%%%%%%%%%%%%%%%%%%%%%%%%%%%%%%%%%%%%%%%%%%%%%%%%%%%%%%%

For equality constraints, a typical SQP step solves

\[
\boxed{
\begin{aligned}
\min_{\mathbf{p}}
\quad&
\nabla f(\mathbf{x}_k)^T\mathbf{p}
+
\frac12
\mathbf{p}^TB_k\mathbf{p}
\\
\text{subject to}
\quad&
h(\mathbf{x}_k)
+
J_h(\mathbf{x}_k)\mathbf{p}
=
0.
\end{aligned}
}
\]

Here \(B_k\) approximates the Hessian of the Lagrangian.

%%%%%%%%%%%%%%%%%%%%%%%%%%%%%%%%%%%%%%%%%%%%%%%%%%%%%%%%%%%%
\subsection{SQP with Inequality Constraints}
\label{subsec:sqp-inequality}
%%%%%%%%%%%%%%%%%%%%%%%%%%%%%%%%%%%%%%%%%%%%%%%%%%%%%%%%%%%%

For nonlinear inequalities,

\[
g_i(\mathbf{x})\leq0,
\]

the SQP subproblem uses linearized constraints:

\[
\boxed{
g_i(\mathbf{x}_k)
+
\nabla g_i(\mathbf{x}_k)^T\mathbf{p}
\leq0.
}
\]

Thus each nonlinear iteration solves a quadratic program with locally
linearized constraints.

%%%%%%%%%%%%%%%%%%%%%%%%%%%%%%%%%%%%%%%%%%%%%%%%%%%%%%%%%%%%
\subsection{Why SQP Is Powerful}
\label{subsec:why-sqp}
%%%%%%%%%%%%%%%%%%%%%%%%%%%%%%%%%%%%%%%%%%%%%%%%%%%%%%%%%%%%

Near a regular solution, SQP closely resembles Newton's method applied to
the KKT equations.

With accurate Hessian information and suitable assumptions, SQP can have
very fast local convergence.

It is one of the standard methods for smooth constrained nonlinear
optimization.

%%%%%%%%%%%%%%%%%%%%%%%%%%%%%%%%%%%%%%%%%%%%%%%%%%%%%%%%%%%%
\subsection{Quasi-Newton SQP}
\label{subsec:quasi-newton-sqp}
%%%%%%%%%%%%%%%%%%%%%%%%%%%%%%%%%%%%%%%%%%%%%%%%%%%%%%%%%%%%

Computing

\[
\nabla_{\mathbf{x}\mathbf{x}}^2
\mathcal{L}
\]

exactly can be expensive.

Instead, \(B_k\) may be updated using quasi-Newton formulas such as BFGS.

This reduces derivative cost while retaining curvature information.

%%%%%%%%%%%%%%%%%%%%%%%%%%%%%%%%%%%%%%%%%%%%%%%%%%%%%%%%%%%%
\subsection{Line Search in Constrained Optimization}
\label{subsec:nlp-line-search}
%%%%%%%%%%%%%%%%%%%%%%%%%%%%%%%%%%%%%%%%%%%%%%%%%%%%%%%%%%%%

After computing a search direction \(\mathbf{p}_k\), a line-search method
chooses

\[
\alpha_k>0
\]

and updates

\[
\boxed{
\mathbf{x}_{k+1}
=
\mathbf{x}_k
+
\alpha_k\mathbf{p}_k.
}
\]

The step should improve some measure of both objective quality and
constraint satisfaction.

%%%%%%%%%%%%%%%%%%%%%%%%%%%%%%%%%%%%%%%%%%%%%%%%%%%%%%%%%%%%
\subsection{Merit Functions}
\label{subsec:nlp-merit-functions}
%%%%%%%%%%%%%%%%%%%%%%%%%%%%%%%%%%%%%%%%%%%%%%%%%%%%%%%%%%%%

A merit function combines objective and constraint violation into one
scalar quantity.

For example,

\[
\boxed{
\Phi_\rho(\mathbf{x})
=
f(\mathbf{x})
+
\rho
\|
h(\mathbf{x})
\|_1.
}
\]

Line search can then require sufficient decrease in \(\Phi_\rho\).

%%%%%%%%%%%%%%%%%%%%%%%%%%%%%%%%%%%%%%%%%%%%%%%%%%%%%%%%%%%%
\subsection{Why Merit Functions Are Needed}
\label{subsec:why-merit-functions}
%%%%%%%%%%%%%%%%%%%%%%%%%%%%%%%%%%%%%%%%%%%%%%%%%%%%%%%%%%%%

A step can reduce the objective while worsening feasibility.

Another step can improve feasibility while increasing the objective.

A merit function gives one scalar quantity for balancing these competing
goals during globalization.

%%%%%%%%%%%%%%%%%%%%%%%%%%%%%%%%%%%%%%%%%%%%%%%%%%%%%%%%%%%%
\subsection{Filter Methods Preview}
\label{subsec:filter-methods-preview}
%%%%%%%%%%%%%%%%%%%%%%%%%%%%%%%%%%%%%%%%%%%%%%%%%%%%%%%%%%%%

Filter methods avoid combining objective and feasibility into one fixed
merit function.

Instead, a trial point can be accepted if it sufficiently improves either:

\[
\boxed{
\text{objective value}
}
\]

or

\[
\boxed{
\text{constraint violation}.
}
\]

This creates a two-criterion acceptance strategy.

%%%%%%%%%%%%%%%%%%%%%%%%%%%%%%%%%%%%%%%%%%%%%%%%%%%%%%%%%%%%
\subsection{Trust-Region Constrained Methods}
\label{subsec:nlp-trust-region}
%%%%%%%%%%%%%%%%%%%%%%%%%%%%%%%%%%%%%%%%%%%%%%%%%%%%%%%%%%%%

Trust-region methods solve a local constrained model subject to

\[
\boxed{
\|\mathbf{p}\|
\leq
\Delta_k.
}
\]

The step is restricted to a region where the local approximation is
believed to be reliable.

The radius is adjusted based on agreement between model prediction and
actual improvement.

%%%%%%%%%%%%%%%%%%%%%%%%%%%%%%%%%%%%%%%%%%%%%%%%%%%%%%%%%%%%
\subsection{Feasibility Restoration}
\label{subsec:feasibility-restoration}
%%%%%%%%%%%%%%%%%%%%%%%%%%%%%%%%%%%%%%%%%%%%%%%%%%%%%%%%%%%%

Some nonlinear algorithms temporarily focus primarily on reducing
constraint violation.

A restoration phase may approximately solve

\[
\boxed{
\min_{\mathbf{x}}
\|
h(\mathbf{x})
\|
}
\]

and suitable measures of violated inequalities.

Once feasibility improves, objective optimization resumes.

%%%%%%%%%%%%%%%%%%%%%%%%%%%%%%%%%%%%%%%%%%%%%%%%%%%%%%%%%%%%
\subsection{Nonlinear Interior-Point Methods}
\label{subsec:nonlinear-interior-point}
%%%%%%%%%%%%%%%%%%%%%%%%%%%%%%%%%%%%%%%%%%%%%%%%%%%%%%%%%%%%

Interior-point methods for nonlinear programming solve barrier-modified
KKT systems.

Introduce slack variables

\[
s_i>0
\]

so that

\[
\boxed{
g_i(\mathbf{x})+s_i=0.
}
\]

Perturbed complementary slackness takes form

\[
\boxed{
\mu_i s_i
=
\tau,
}
\]

where

\[
\tau>0
\]

is driven toward zero.

%%%%%%%%%%%%%%%%%%%%%%%%%%%%%%%%%%%%%%%%%%%%%%%%%%%%%%%%%%%%
\subsection{Primal-Dual KKT System}
\label{subsec:nlp-primal-dual-kkt}
%%%%%%%%%%%%%%%%%%%%%%%%%%%%%%%%%%%%%%%%%%%%%%%%%%%%%%%%%%%%

A primal-dual nonlinear interior-point method seeks approximate
solutions of

\[
\boxed{
\nabla f
+
J_g^T\boldsymbol{\mu}
+
J_h^T\boldsymbol{\lambda}
=
0,
}
\]

\[
\boxed{
g(\mathbf{x})+\mathbf{s}=0,
}
\]

\[
\boxed{
h(\mathbf{x})=0,
}
\]

and

\[
\boxed{
S\boldsymbol{\mu}
=
\tau\mathbf{1}.
}
\]

Here

\[
S
=
\operatorname{diag}(s_1,\ldots,s_m).
\]

%%%%%%%%%%%%%%%%%%%%%%%%%%%%%%%%%%%%%%%%%%%%%%%%%%%%%%%%%%%%
\subsection{Newton's Method for the Primal-Dual System}
\label{subsec:nlp-primal-dual-newton}
%%%%%%%%%%%%%%%%%%%%%%%%%%%%%%%%%%%%%%%%%%%%%%%%%%%%%%%%%%%%

The nonlinear primal-dual equations are linearized around the current
iterate.

A Newton system is solved for corrections to:

\[
\boxed{
\mathbf{x},
}
\]

\[
\boxed{
\boldsymbol{\lambda},
}
\]

\[
\boxed{
\boldsymbol{\mu},
}
\]

and

\[
\boxed{
\mathbf{s}.
}
\]

Large practical NLP solvers rely heavily on sparse linear algebra for this
step.

%%%%%%%%%%%%%%%%%%%%%%%%%%%%%%%%%%%%%%%%%%%%%%%%%%%%%%%%%%%%
\subsection{Active-Set Methods}
\label{subsec:nlp-active-set-methods}
%%%%%%%%%%%%%%%%%%%%%%%%%%%%%%%%%%%%%%%%%%%%%%%%%%%%%%%%%%%%

An active-set method attempts to identify which inequalities will be
binding at the solution.

Once an active set is guessed, the problem is treated locally like an
equality-constrained problem.

The algorithm then adds or removes constraints from the active set as
needed.

%%%%%%%%%%%%%%%%%%%%%%%%%%%%%%%%%%%%%%%%%%%%%%%%%%%%%%%%%%%%
\subsection{Projected Gradient Methods Preview}
\label{subsec:projected-gradient-preview}
%%%%%%%%%%%%%%%%%%%%%%%%%%%%%%%%%%%%%%%%%%%%%%%%%%%%%%%%%%%%

For a simple feasible set \(C\), projected gradient descent uses

\[
\boxed{
\mathbf{x}_{k+1}
=
P_C
\left[
\mathbf{x}_k
-
\alpha_k
\nabla f(\mathbf{x}_k)
\right],
}
\]

where

\[
P_C
\]

denotes projection onto \(C\).

This is especially useful when projection is inexpensive.

%%%%%%%%%%%%%%%%%%%%%%%%%%%%%%%%%%%%%%%%%%%%%%%%%%%%%%%%%%%%
\subsection{Bound-Constrained Optimization}
\label{subsec:bound-constrained-optimization}
%%%%%%%%%%%%%%%%%%%%%%%%%%%%%%%%%%%%%%%%%%%%%%%%%%%%%%%%%%%%

A common nonlinear problem has simple bounds:

\[
\boxed{
\ell_i
\leq
x_i
\leq
u_i.
}
\]

These constraints arise naturally from:

\[
\boxed{
\text{physical limits},
}
\]

\[
\boxed{
\text{probabilities},
}
\]

and

\[
\boxed{
\text{design ranges}.
}
\]

Specialized algorithms exploit this simple structure.

%%%%%%%%%%%%%%%%%%%%%%%%%%%%%%%%%%%%%%%%%%%%%%%%%%%%%%%%%%%%
\subsection{Variable Transformations for Bounds}
\label{subsec:nlp-variable-transformations}
%%%%%%%%%%%%%%%%%%%%%%%%%%%%%%%%%%%%%%%%%%%%%%%%%%%%%%%%%%%%

Sometimes constraints are removed through reparameterization.

For a positive variable,

\[
x>0,
\]

one may write

\[
\boxed{
x=e^z.
}
\]

For

\[
0<x<1,
\]

one may use the logistic transformation

\[
\boxed{
x
=
\frac{
1
}{
1+e^{-z}
}.
}
\]

Such transformations trade explicit constraints for nonlinear geometry.

%%%%%%%%%%%%%%%%%%%%%%%%%%%%%%%%%%%%%%%%%%%%%%%%%%%%%%%%%%%%
\subsection{Nonlinear Least Squares}
\label{subsec:nonlinear-least-squares}
%%%%%%%%%%%%%%%%%%%%%%%%%%%%%%%%%%%%%%%%%%%%%%%%%%%%%%%%%%%%

Suppose residuals are

\[
r_i(\mathbf{x}),
\qquad
i=1,\ldots,m.
\]

Nonlinear least squares solves

\[
\boxed{
\min_{\mathbf{x}}
\frac12
\sum_{i=1}^{m}
r_i(\mathbf{x})^2.
}
\]

In vector form,

\[
\boxed{
\min_{\mathbf{x}}
\frac12
\|
\mathbf{r}(\mathbf{x})
\|_2^2.
}
\]

%%%%%%%%%%%%%%%%%%%%%%%%%%%%%%%%%%%%%%%%%%%%%%%%%%%%%%%%%%%%
\subsection{Gradient of Nonlinear Least Squares}
\label{subsec:nonlinear-least-squares-gradient}
%%%%%%%%%%%%%%%%%%%%%%%%%%%%%%%%%%%%%%%%%%%%%%%%%%%%%%%%%%%%

Let

\[
J(\mathbf{x})
\]

be the Jacobian of the residual vector.

Then

\[
\boxed{
\nabla f(\mathbf{x})
=
J(\mathbf{x})^T
\mathbf{r}(\mathbf{x}).
}
\]

The Hessian is

\[
\boxed{
\nabla^2f
=
J^TJ
+
\sum_i
r_i
\nabla^2r_i.
}
\]

%%%%%%%%%%%%%%%%%%%%%%%%%%%%%%%%%%%%%%%%%%%%%%%%%%%%%%%%%%%%
\subsection{Gauss--Newton Method}
\label{subsec:gauss-newton-method}
%%%%%%%%%%%%%%%%%%%%%%%%%%%%%%%%%%%%%%%%%%%%%%%%%%%%%%%%%%%%

If residuals are small or nearly linear, the second Hessian term may be
neglected.

The Gauss--Newton step solves

\[
\boxed{
J^TJ\mathbf{p}
=
-
J^T\mathbf{r}.
}
\]

This is often effective for parameter estimation.

%%%%%%%%%%%%%%%%%%%%%%%%%%%%%%%%%%%%%%%%%%%%%%%%%%%%%%%%%%%%
\subsection{Levenberg--Marquardt Preview}
\label{subsec:levenberg-marquardt}
%%%%%%%%%%%%%%%%%%%%%%%%%%%%%%%%%%%%%%%%%%%%%%%%%%%%%%%%%%%%

Levenberg--Marquardt modifies the Gauss--Newton system:

\[
\boxed{
(
J^TJ+\lambda I
)
\mathbf{p}
=
-
J^T\mathbf{r}.
}
\]

The parameter \(\lambda>0\) regularizes the step.

Large \(\lambda\) makes the method more gradient-like, while small
\(\lambda\) approaches Gauss--Newton behavior.

%%%%%%%%%%%%%%%%%%%%%%%%%%%%%%%%%%%%%%%%%%%%%%%%%%%%%%%%%%%%
\subsection{Globalization}
\label{subsec:nlp-globalization}
%%%%%%%%%%%%%%%%%%%%%%%%%%%%%%%%%%%%%%%%%%%%%%%%%%%%%%%%%%%%

Newton and SQP methods can converge rapidly near a solution but may fail
when started far away.

Globalization strategies improve robustness.

Common strategies include:

\[
\boxed{
\text{line search},
}
\]

\[
\boxed{
\text{trust regions},
}
\]

\[
\boxed{
\text{merit functions},
}
\]

and

\[
\boxed{
\text{filters}.
}
\]

Here \emph{globalization} refers to convergence from a broader set of
starting points, not proof of a global optimum.

%%%%%%%%%%%%%%%%%%%%%%%%%%%%%%%%%%%%%%%%%%%%%%%%%%%%%%%%%%%%
\subsection{Local Solvers}
\label{subsec:nlp-local-solvers}
%%%%%%%%%%%%%%%%%%%%%%%%%%%%%%%%%%%%%%%%%%%%%%%%%%%%%%%%%%%%

Most classical smooth NLP algorithms are local methods.

They attempt to find a nearby KKT point.

The solution can depend strongly on:

\[
\boxed{
\text{initial guess},
}
\]

\[
\boxed{
\text{scaling},
}
\]

and

\[
\boxed{
\text{algorithmic parameters}.
}
\]

%%%%%%%%%%%%%%%%%%%%%%%%%%%%%%%%%%%%%%%%%%%%%%%%%%%%%%%%%%%%
\subsection{Multistart Methods}
\label{subsec:nlp-multistart}
%%%%%%%%%%%%%%%%%%%%%%%%%%%%%%%%%%%%%%%%%%%%%%%%%%%%%%%%%%%%

A simple strategy for nonconvex problems is to run a local solver from
many initial points.

If several runs converge to different local minima, compare their
objective values.

Multistart does not generally prove global optimality, but it can improve
the chance of locating better solutions.

%%%%%%%%%%%%%%%%%%%%%%%%%%%%%%%%%%%%%%%%%%%%%%%%%%%%%%%%%%%%
\subsection{Global Nonlinear Optimization Preview}
\label{subsec:global-nonlinear-optimization-preview}
%%%%%%%%%%%%%%%%%%%%%%%%%%%%%%%%%%%%%%%%%%%%%%%%%%%%%%%%%%%%

Global optimization seeks a provably best solution over the entire
feasible set.

Methods can use:

\[
\boxed{
\text{branch-and-bound},
}
\]

\[
\boxed{
\text{convex relaxations},
}
\]

\[
\boxed{
\text{interval methods},
}
\]

or

\[
\boxed{
\text{spatial decomposition}.
}
\]

Global nonlinear optimization can be substantially harder than local
optimization.

%%%%%%%%%%%%%%%%%%%%%%%%%%%%%%%%%%%%%%%%%%%%%%%%%%%%%%%%%%%%
\subsection{Scaling}
\label{subsec:nlp-scaling}
%%%%%%%%%%%%%%%%%%%%%%%%%%%%%%%%%%%%%%%%%%%%%%%%%%%%%%%%%%%%

Suppose one variable naturally has magnitude

\[
10^{-6}
\]

and another has magnitude

\[
10^6.
\]

Then gradients and Hessians can be poorly balanced.

Variable and constraint scaling can improve:

\[
\boxed{
\text{conditioning},
}
\]

\[
\boxed{
\text{Newton steps},
}
\]

and

\[
\boxed{
\text{termination tests}.
}
\]

%%%%%%%%%%%%%%%%%%%%%%%%%%%%%%%%%%%%%%%%%%%%%%%%%%%%%%%%%%%%
\subsection{Constraint Scaling}
\label{subsec:nlp-constraint-scaling}
%%%%%%%%%%%%%%%%%%%%%%%%%%%%%%%%%%%%%%%%%%%%%%%%%%%%%%%%%%%%

Suppose two constraints are

\[
g_1(\mathbf{x})\approx10^8,
\]

and

\[
g_2(\mathbf{x})\approx10^{-5}.
\]

A solver may treat residuals unevenly.

Rescaling constraints so their typical magnitudes are comparable often
improves numerical behavior.

%%%%%%%%%%%%%%%%%%%%%%%%%%%%%%%%%%%%%%%%%%%%%%%%%%%%%%%%%%%%
\subsection{Sparse Nonlinear Programs}
\label{subsec:sparse-nonlinear-programs}
%%%%%%%%%%%%%%%%%%%%%%%%%%%%%%%%%%%%%%%%%%%%%%%%%%%%%%%%%%%%

Large engineering NLPs often have sparse derivative matrices.

That is, most entries of

\[
\boxed{
J_g,
}
\]

\[
\boxed{
J_h,
}
\]

and

\[
\boxed{
\nabla^2\mathcal{L}
}
\]

are zero.

Sparse linear algebra can dramatically reduce memory and computational
cost.

%%%%%%%%%%%%%%%%%%%%%%%%%%%%%%%%%%%%%%%%%%%%%%%%%%%%%%%%%%%%
\subsection{Automatic Differentiation}
\label{subsec:nlp-automatic-differentiation}
%%%%%%%%%%%%%%%%%%%%%%%%%%%%%%%%%%%%%%%%%%%%%%%%%%%%%%%%%%%%

Optimization methods often require accurate derivatives.

Automatic differentiation evaluates derivatives through the computational
graph of the objective and constraints.

It can provide derivatives accurate to floating-point precision without
symbolic expression growth or finite-difference truncation error.

%%%%%%%%%%%%%%%%%%%%%%%%%%%%%%%%%%%%%%%%%%%%%%%%%%%%%%%%%%%%
\subsection{Finite-Difference Derivatives}
\label{subsec:nlp-finite-difference-derivatives}
%%%%%%%%%%%%%%%%%%%%%%%%%%%%%%%%%%%%%%%%%%%%%%%%%%%%%%%%%%%%

When analytical derivatives are unavailable, approximate

\[
\boxed{
\frac{\partial f}{\partial x_i}
\approx
\frac{
f(\mathbf{x}+h\mathbf{e}_i)
-
f(\mathbf{x})
}{
h
}.
}
\]

However, choosing \(h\) involves a tradeoff between:

\[
\boxed{
\text{truncation error}
}
\]

and

\[
\boxed{
\text{floating-point cancellation}.
}
\]

Poor derivative estimates can seriously degrade nonlinear solvers.

%%%%%%%%%%%%%%%%%%%%%%%%%%%%%%%%%%%%%%%%%%%%%%%%%%%%%%%%%%%%
\subsection{Stopping Criteria}
\label{subsec:nlp-stopping-criteria}
%%%%%%%%%%%%%%%%%%%%%%%%%%%%%%%%%%%%%%%%%%%%%%%%%%%%%%%%%%%%

For constrained nonlinear optimization, stopping criteria should examine
several residuals.

Important measures include:

\[
\boxed{
\|\nabla_{\mathbf{x}}\mathcal{L}\|,
}
\]

\[
\boxed{
\|h(\mathbf{x})\|,
}
\]

\[
\boxed{
\max_i
\{
g_i(\mathbf{x}),0
\},
}
\]

and

\[
\boxed{
|\mu_i g_i(\mathbf{x})|.
}
\]

A small change in the objective alone is not enough.

%%%%%%%%%%%%%%%%%%%%%%%%%%%%%%%%%%%%%%%%%%%%%%%%%%%%%%%%%%%%
\subsection{KKT Residual}
\label{subsec:nlp-kkt-residual}
%%%%%%%%%%%%%%%%%%%%%%%%%%%%%%%%%%%%%%%%%%%%%%%%%%%%%%%%%%%%

A solver may combine stationarity, feasibility, and complementarity into a
KKT residual.

Schematically,

\[
\boxed{
R_{\mathrm{KKT}}
=
\begin{pmatrix}
\nabla_{\mathbf{x}}\mathcal{L}
\\
h(\mathbf{x})
\\
\text{inequality residual}
\\
\text{complementarity residual}
\end{pmatrix}.
}
\]

A small norm

\[
\boxed{
\|R_{\mathrm{KKT}}\|
}
\]

indicates approximate first-order optimality.

%%%%%%%%%%%%%%%%%%%%%%%%%%%%%%%%%%%%%%%%%%%%%%%%%%%%%%%%%%%%
\subsection{Sensitivity Analysis}
\label{subsec:nlp-sensitivity-analysis}
%%%%%%%%%%%%%%%%%%%%%%%%%%%%%%%%%%%%%%%%%%%%%%%%%%%%%%%%%%%%

Suppose the optimization problem depends on a parameter \(p\).

The optimal solution

\[
\mathbf{x}^*(p)
\]

and multipliers

\[
\boldsymbol{\lambda}^*(p),
\qquad
\boldsymbol{\mu}^*(p)
\]

may change as \(p\) changes.

Under suitable regularity conditions, differentiating the KKT system can
provide local sensitivity information.

%%%%%%%%%%%%%%%%%%%%%%%%%%%%%%%%%%%%%%%%%%%%%%%%%%%%%%%%%%%%
\subsection{Shadow Prices in Nonlinear Programming}
\label{subsec:nlp-shadow-prices}
%%%%%%%%%%%%%%%%%%%%%%%%%%%%%%%%%%%%%%%%%%%%%%%%%%%%%%%%%%%%

Lagrange multipliers can still have marginal-value interpretations.

If a constraint is parameterized by

\[
h(\mathbf{x})=b,
\]

then the associated multiplier may describe the local change in optimal
value caused by changing \(b\).

The exact sign depends on the formulation.

Unlike linear programming, this sensitivity may vary continuously with
the operating point.

%%%%%%%%%%%%%%%%%%%%%%%%%%%%%%%%%%%%%%%%%%%%%%%%%%%%%%%%%%%%
\subsection{Engineering Design Example}
\label{subsec:nlp-engineering-design}
%%%%%%%%%%%%%%%%%%%%%%%%%%%%%%%%%%%%%%%%%%%%%%%%%%%%%%%%%%%%

Suppose a cylindrical container must hold volume \(V\).

Let:

\[
r
=
\text{radius},
\]

and

\[
h
=
\text{height}.
\]

Volume imposes

\[
\boxed{
\pi r^2h=V.
}
\]

Surface area is

\[
\boxed{
A(r,h)
=
2\pi r^2
+
2\pi rh.
}
\]

The design problem is

\[
\boxed{
\begin{aligned}
\min_{r,h}
\quad&
2\pi r^2+2\pi rh
\\
\text{subject to}
\quad&
\pi r^2h=V,
\\
&
r,h>0.
\end{aligned}
}
\]

This is a nonlinear equality-constrained optimization problem.

%%%%%%%%%%%%%%%%%%%%%%%%%%%%%%%%%%%%%%%%%%%%%%%%%%%%%%%%%%%%
\subsection{Worked Example: Optimal Cylinder Proportion}
\label{subsec:worked-optimal-cylinder}
%%%%%%%%%%%%%%%%%%%%%%%%%%%%%%%%%%%%%%%%%%%%%%%%%%%%%%%%%%%%

\begin{workedexamplebox}{Closed Cylinder with Fixed Volume}

From

\[
\pi r^2h=V,
\]

solve for

\[
h
=
\frac{
V
}{
\pi r^2
}.
\]

Substitute into surface area:

\[
A(r)
=
2\pi r^2
+
2\pi r
\left(
\frac{
V
}{
\pi r^2
}
\right).
\]

Thus,

\[
A(r)
=
2\pi r^2
+
\frac{
2V
}{
r
}.
\]

Differentiate:

\[
A'(r)
=
4\pi r
-
\frac{
2V
}{
r^2
}.
\]

Set equal to zero:

\[
4\pi r^3
=
2V.
\]

Therefore,

\[
r^3
=
\frac{
V
}{
2\pi
}.
\]

Since

\[
h
=
\frac{
V
}{
\pi r^2
},
\]

using

\[
V=2\pi r^3
\]

gives

\[
h=2r.
\]

Hence,

\begin{answerbox}
\[
\boxed{
h=2r.
}
\]

The optimal cylinder height equals its diameter.
\end{answerbox}

\end{workedexamplebox}

%%%%%%%%%%%%%%%%%%%%%%%%%%%%%%%%%%%%%%%%%%%%%%%%%%%%%%%%%%%%
\subsection{Parameter Estimation}
\label{subsec:nlp-parameter-estimation}
%%%%%%%%%%%%%%%%%%%%%%%%%%%%%%%%%%%%%%%%%%%%%%%%%%%%%%%%%%%%

Suppose a nonlinear model predicts

\[
y_i
\approx
m(t_i,\boldsymbol{\theta}).
\]

Parameter estimation can be formulated as

\[
\boxed{
\min_{\boldsymbol{\theta}}
\sum_{i=1}^{n}
\left[
y_i
-
m(t_i,\boldsymbol{\theta})
\right]^2.
}
\]

This is nonlinear least squares when \(m\) depends nonlinearly on the
parameters.

%%%%%%%%%%%%%%%%%%%%%%%%%%%%%%%%%%%%%%%%%%%%%%%%%%%%%%%%%%%%
\subsection{Maximum Likelihood as Nonlinear Optimization}
\label{subsec:nlp-maximum-likelihood}
%%%%%%%%%%%%%%%%%%%%%%%%%%%%%%%%%%%%%%%%%%%%%%%%%%%%%%%%%%%%

For likelihood

\[
L(\boldsymbol{\theta}),
\]

maximum likelihood solves

\[
\boxed{
\max_{\boldsymbol{\theta}}
L(\boldsymbol{\theta}).
}
\]

Equivalently,

\[
\boxed{
\min_{\boldsymbol{\theta}}
-
\log
L(\boldsymbol{\theta}).
}
\]

Constraints may arise from parameter domains such as:

\[
\boxed{
\sigma>0,
}
\]

or

\[
\boxed{
p\in[0,1].
}
\]

%%%%%%%%%%%%%%%%%%%%%%%%%%%%%%%%%%%%%%%%%%%%%%%%%%%%%%%%%%%%
\subsection{Economic Optimization}
\label{subsec:nlp-economic-optimization}
%%%%%%%%%%%%%%%%%%%%%%%%%%%%%%%%%%%%%%%%%%%%%%%%%%%%%%%%%%%%

Consumer and producer models often involve nonlinear utility, production,
or cost functions.

For example,

\[
\boxed{
\max_{x,y}
x^\alpha y^{1-\alpha}
}
\]

subject to

\[
\boxed{
p_xx+p_yy\leq I.
}
\]

This is a nonlinear optimization problem with a linear budget constraint.

%%%%%%%%%%%%%%%%%%%%%%%%%%%%%%%%%%%%%%%%%%%%%%%%%%%%%%%%%%%%
\subsection{Nonlinear Programming Workflow}
\label{subsec:nlp-workflow}
%%%%%%%%%%%%%%%%%%%%%%%%%%%%%%%%%%%%%%%%%%%%%%%%%%%%%%%%%%%%

A practical workflow is:

\begin{enumerate}
    \item define the decision variables;
    \item write the nonlinear objective;
    \item state equality and inequality constraints;
    \item specify bounds and domains;
    \item determine whether the problem is smooth;
    \item determine whether the problem is convex;
    \item find or construct a feasible starting point;
    \item identify active constraints near candidate solutions;
    \item derive first-order KKT conditions;
    \item check constraint qualifications;
    \item examine second-order conditions when possible;
    \item select an algorithm appropriate to the problem structure;
    \item scale variables and constraints;
    \item solve the numerical problem;
    \item inspect KKT residuals;
    \item test multiple initial points if nonconvexity is present;
    \item perform sensitivity analysis;
    \item interpret the result in the original application.
\end{enumerate}

%%%%%%%%%%%%%%%%%%%%%%%%%%%%%%%%%%%%%%%%%%%%%%%%%%%%%%%%%%%%
\subsection{Choosing a Nonlinear Programming Method}
\label{subsec:choosing-nlp-method}
%%%%%%%%%%%%%%%%%%%%%%%%%%%%%%%%%%%%%%%%%%%%%%%%%%%%%%%%%%%%

For smooth equality constraints, consider:

\[
\boxed{
\text{Newton or SQP methods}.
}
\]

For many smooth inequalities, consider:

\[
\boxed{
\text{interior-point methods}.
}
\]

For difficult feasibility structure, consider:

\[
\boxed{
\text{augmented Lagrangian methods}.
}
\]

For simple bounds, consider:

\[
\boxed{
\text{projected or bound-constrained methods}.
}
\]

For nonlinear least squares, consider:

\[
\boxed{
\text{Gauss--Newton or Levenberg--Marquardt}.
}
\]

For nonconvex global optimization, specialized global methods may be
necessary.

%%%%%%%%%%%%%%%%%%%%%%%%%%%%%%%%%%%%%%%%%%%%%%%%%%%%%%%%%%%%
\subsection{Common Errors}
\label{subsec:nonlinear-programming-common-errors}
%%%%%%%%%%%%%%%%%%%%%%%%%%%%%%%%%%%%%%%%%%%%%%%%%%%%%%%%%%%%

\subsubsection{Assuming KKT Conditions Guarantee a Minimum}

For a nonconvex problem, a KKT point need not be a minimum.

Second-order conditions or global reasoning may be required.

%%%%%%%%%%%%%%%%%%%%%%%%%%%%%%%%%%%%%%%%%%%%%%%%%%%%%%%%%%%%
\subsubsection{Ignoring Constraint Qualifications}

KKT conditions may fail to be necessary when constraint gradients are
degenerate.

Always examine regularity assumptions.

%%%%%%%%%%%%%%%%%%%%%%%%%%%%%%%%%%%%%%%%%%%%%%%%%%%%%%%%%%%%
\subsubsection{Checking the Full Hessian Instead of Feasible Directions}

For constrained optimization, curvature matters along feasible tangent or
critical directions.

The full Hessian need not be positive definite.

%%%%%%%%%%%%%%%%%%%%%%%%%%%%%%%%%%%%%%%%%%%%%%%%%%%%%%%%%%%%
\subsubsection{Using the Hessian of \(f\) Instead of the Hessian of the Lagrangian}

Second-order constrained conditions generally use

\[
\boxed{
\nabla_{\mathbf{x}\mathbf{x}}^2
\mathcal{L},
}
\]

not merely

\[
\nabla^2f.
\]

%%%%%%%%%%%%%%%%%%%%%%%%%%%%%%%%%%%%%%%%%%%%%%%%%%%%%%%%%%%%
\subsubsection{Forgetting Complementary Slackness}

For every inequality,

\[
\boxed{
\mu_i g_i(x)=0.
}
\]

An inactive constraint must have zero multiplier.

%%%%%%%%%%%%%%%%%%%%%%%%%%%%%%%%%%%%%%%%%%%%%%%%%%%%%%%%%%%%
\subsubsection{Using the Wrong Multiplier Sign}

For

\[
g_i(x)\leq0,
\]

the standard minimization KKT convention requires

\[
\boxed{
\mu_i\geq0.
}
\]

Changing the inequality orientation changes the sign convention.

%%%%%%%%%%%%%%%%%%%%%%%%%%%%%%%%%%%%%%%%%%%%%%%%%%%%%%%%%%%%
\subsubsection{Assuming Penalty Methods Enforce Exact Feasibility at Finite \(\rho\)}

Quadratic penalties generally require increasing \(\rho\) toward large
values.

At finite penalty, constraint violation is usually nonzero.

%%%%%%%%%%%%%%%%%%%%%%%%%%%%%%%%%%%%%%%%%%%%%%%%%%%%%%%%%%%%
\subsubsection{Taking Penalty Parameters Too Large Too Quickly}

Very large \(\rho\) can create ill-conditioning.

Augmented Lagrangian methods help reduce this problem.

%%%%%%%%%%%%%%%%%%%%%%%%%%%%%%%%%%%%%%%%%%%%%%%%%%%%%%%%%%%%
\subsubsection{Using a Log Barrier Outside the Strict Interior}

The logarithmic barrier

\[
-\log[-g_i(x)]
\]

requires

\[
\boxed{
g_i(x)<0.
}
\]

It is undefined on or outside the boundary.

%%%%%%%%%%%%%%%%%%%%%%%%%%%%%%%%%%%%%%%%%%%%%%%%%%%%%%%%%%%%
\subsubsection{Confusing Globalization with Global Optimization}

Line search and trust regions improve convergence reliability.

They do not generally guarantee the global optimum of a nonconvex
problem.

%%%%%%%%%%%%%%%%%%%%%%%%%%%%%%%%%%%%%%%%%%%%%%%%%%%%%%%%%%%%
\subsubsection{Ignoring Initial Guess Dependence}

Local nonlinear solvers can converge to different KKT points from
different starting values.

This is especially important for nonconvex problems.

%%%%%%%%%%%%%%%%%%%%%%%%%%%%%%%%%%%%%%%%%%%%%%%%%%%%%%%%%%%%
\subsubsection{Ignoring Scaling}

Poorly scaled variables and constraints can cause bad Newton steps and
misleading termination tests.

%%%%%%%%%%%%%%%%%%%%%%%%%%%%%%%%%%%%%%%%%%%%%%%%%%%%%%%%%%%%
\subsubsection{Using Noisy Finite Differences Without Checking Accuracy}

Derivative errors can make an otherwise good optimization algorithm fail.

Analytical or automatic derivatives are preferable when available.

%%%%%%%%%%%%%%%%%%%%%%%%%%%%%%%%%%%%%%%%%%%%%%%%%%%%%%%%%%%%
\subsubsection{Stopping Because the Objective Barely Changes}

A flat objective sequence does not prove feasibility or stationarity.

KKT residuals should also be checked.

%%%%%%%%%%%%%%%%%%%%%%%%%%%%%%%%%%%%%%%%%%%%%%%%%%%%%%%%%%%%
\subsection{Applications}
\label{subsec:nonlinear-programming-applications}
%%%%%%%%%%%%%%%%%%%%%%%%%%%%%%%%%%%%%%%%%%%%%%%%%%%%%%%%%%%%

\begin{applicationbox}

\textbf{Engineering Design.}
Structures, aircraft, mechanical components, and energy systems are often
optimized subject to nonlinear physical constraints.

\medskip

\textbf{Parameter Estimation.}
Nonlinear least-squares problems arise when model predictions depend
nonlinearly on unknown parameters.

\medskip

\textbf{Statistics.}
Maximum-likelihood and penalized estimation frequently require nonlinear
optimization.

\medskip

\textbf{Machine Learning.}
Training many models requires minimizing large nonlinear and often
nonconvex loss functions.

\medskip

\textbf{Economics.}
Utility, production, equilibrium, and resource-allocation models often
contain nonlinear objectives and constraints.

\medskip

\textbf{Finance.}
Risk measures, transaction costs, and nonlinear market constraints can
produce nonlinear portfolio models.

\medskip

\textbf{Optimal Control.}
Discretized trajectory optimization typically produces large sparse NLPs.

\medskip

\textbf{Scientific Computing.}
Inverse problems, model calibration, and experimental design frequently
lead to constrained nonlinear optimization.

\end{applicationbox}

%%%%%%%%%%%%%%%%%%%%%%%%%%%%%%%%%%%%%%%%%%%%%%%%%%%%%%%%%%%%
\subsection{Exercises}
\label{subsec:nonlinear-programming-exercises}
%%%%%%%%%%%%%%%%%%%%%%%%%%%%%%%%%%%%%%%%%%%%%%%%%%%%%%%%%%%%

\begin{exercise}[1]
Define a nonlinear programming problem.
\end{exercise}

\begin{exercise}[1]
Give an example of a nonlinear objective with linear constraints.
\end{exercise}

\begin{exercise}[1]
Give an example of a linear objective with nonlinear constraints.
\end{exercise}

\begin{exercise}[1]
Define an active inequality constraint.
\end{exercise}

\begin{exercise}[1]
Define LICQ conceptually.
\end{exercise}

\begin{exercise}[1]
State the KKT conditions.
\end{exercise}

\begin{exercise}[1]
Define the Hessian of the Lagrangian.
\end{exercise}

\begin{exercise}[1]
Define a quadratic penalty function.
\end{exercise}

\begin{exercise}[1]
Define an augmented Lagrangian.
\end{exercise}

\begin{exercise}[1]
Define sequential quadratic programming conceptually.
\end{exercise}

\begin{exercise}[2]
Determine whether

\[
\min_{x,y}
x^2+y
\]

subject to

\[
x+y\leq3
\]

is a nonlinear program.
\end{exercise}

\begin{exercise}[2]
Identify the active constraints at

\[
(x,y)=(1,1)
\]

for

\[
x^2+y^2\leq2,
\]

\[
x+y\leq3.
\]
\end{exercise}

\begin{exercise}[2]
Find the tangent direction condition for

\[
x^2+y^2=1.
\]
\end{exercise}

\begin{exercise}[2]
Write the Lagrangian for

\[
\min_{x,y}
x^2+3y^2
\]

subject to

\[
xy=1.
\]
\end{exercise}

\begin{exercise}[2]
Write the KKT conditions for

\[
\min_x
x^2
\]

subject to

\[
1-x\leq0.
\]
\end{exercise}

\begin{exercise}[3]
Minimize

\[
x^2+y^2
\]

subject to

\[
xy=1
\]

using Lagrange multipliers.
\end{exercise}

\begin{exercise}[3]
Minimize

\[
x+y
\]

subject to

\[
x^2+y^2=4.
\]
\end{exercise}

\begin{exercise}[3]
Find the KKT points of

\[
\min_{x,y}
x^2+y^2
\]

subject to

\[
x+y\geq4.
\]
\end{exercise}

\begin{exercise}[3]
Verify LICQ for the active constraint in the previous problem.
\end{exercise}

\begin{exercise}[3]
Give a constrained optimization problem where LICQ fails.
\end{exercise}

\begin{exercise}[3]
Explain why

\[
\min_x x
\]

subject to

\[
x^2=0
\]

does not satisfy ordinary Lagrange multiplier stationarity.
\end{exercise}

\begin{exercise}[3]
Derive the Hessian of the Lagrangian for

\[
f(x,y)=x^2+y^2
\]

subject to

\[
xy=1.
\]
\end{exercise}

\begin{exercise}[3]
Determine the tangent space to

\[
x+y+z=1.
\]
\end{exercise}

\begin{exercise}[3]
Use the reduced Hessian to analyze an equality-constrained quadratic
problem.
\end{exercise}

\begin{exercise}[3]
Derive the Newton KKT system for equality-constrained optimization.
\end{exercise}

\begin{exercise}[3]
Construct a quadratic penalty formulation for

\[
h(x)=0.
\]
\end{exercise}

\begin{exercise}[3]
Show that a quadratic penalty solution approaches feasibility as
\(\rho\to\infty\) for a simple one-variable example.
\end{exercise}

\begin{exercise}[3]
Construct a logarithmic barrier for

\[
x^2+y^2<1.
\]
\end{exercise}

\begin{exercise}[3]
Write an augmented Lagrangian for two equality constraints.
\end{exercise}

\begin{exercise}[3]
Derive a multiplier update for an augmented Lagrangian method.
\end{exercise}

\begin{exercise}[3]
Construct the SQP subproblem for

\[
\min f(x)
\]

subject to

\[
h(x)=0.
\]
\end{exercise}

\begin{exercise}[3]
Linearize the nonlinear inequality

\[
x^2+y^2-1\leq0
\]

around a given point.
\end{exercise}

\begin{exercise}[3]
Explain why a merit function is useful in constrained line search.
\end{exercise}

\begin{exercise}[3]
Derive the gradient for nonlinear least squares.
\end{exercise}

\begin{exercise}[3]
Derive the exact Hessian of a nonlinear least-squares objective.
\end{exercise}

\begin{exercise}[3]
Derive the Gauss--Newton approximation.
\end{exercise}

\begin{exercise}[3]
Formulate a nonlinear parameter-estimation problem.
\end{exercise}

\begin{exercise}[3]
Formulate a nonlinear engineering-design problem with one equality and
one inequality constraint.
\end{exercise}

\begin{exercise}[3]
Explain why multistart can help on a nonconvex NLP.
\end{exercise}

\begin{exercise}[4]
Prove first-order necessary conditions using tangent cones.
\end{exercise}

\begin{exercise}[4]
Study LICQ, MFCQ, and their relationships.
\end{exercise}

\begin{exercise}[4]
Study second-order necessary conditions for nonlinear programs.
\end{exercise}

\begin{exercise}[4]
Study second-order sufficient conditions using critical cones.
\end{exercise}

\begin{exercise}[4]
Derive reduced-Hessian conditions for equality-constrained problems.
\end{exercise}

\begin{exercise}[4]
Study exact penalty theory.
\end{exercise}

\begin{exercise}[4]
Analyze ill-conditioning in quadratic penalty methods.
\end{exercise}

\begin{exercise}[4]
Study augmented Lagrangian convergence.
\end{exercise}

\begin{exercise}[4]
Derive SQP from Newton's method on the KKT equations.
\end{exercise}

\begin{exercise}[4]
Study superlinear convergence of SQP.
\end{exercise}

\begin{exercise}[4]
Study line-search merit functions.
\end{exercise}

\begin{exercise}[4]
Investigate filter globalization methods.
\end{exercise}

\begin{exercise}[4]
Study trust-region SQP methods.
\end{exercise}

\begin{exercise}[4]
Derive primal-dual interior-point equations for NLP.
\end{exercise}

\begin{exercise}[4]
Study central-path theory for nonlinear programs.
\end{exercise}

\begin{exercise}[4]
Investigate sparse KKT-system factorization.
\end{exercise}

\begin{exercise}[4]
Study automatic differentiation for large NLPs.
\end{exercise}

\begin{exercise}[4]
Analyze Gauss--Newton convergence.
\end{exercise}

\begin{exercise}[4]
Study Levenberg--Marquardt methods.
\end{exercise}

\begin{exercise}[4]
Investigate global optimization methods for nonconvex NLPs.
\end{exercise}

%%%%%%%%%%%%%%%%%%%%%%%%%%%%%%%%%%%%%%%%%%%%%%%%%%%%%%%%%%%%
\subsection{Conceptual Summary}
\label{subsec:nonlinear-programming-summary}
%%%%%%%%%%%%%%%%%%%%%%%%%%%%%%%%%%%%%%%%%%%%%%%%%%%%%%%%%%%%

A nonlinear program has form

\[
\boxed{
\begin{aligned}
\min_{\mathbf{x}}
\quad&
f(\mathbf{x})
\\
\text{subject to}
\quad&
g_i(\mathbf{x})\leq0,
\\
&
h_j(\mathbf{x})=0.
\end{aligned}
}
\]

Its Lagrangian is

\[
\boxed{
\mathcal{L}
=
f
+
\sum_i
\mu_i g_i
+
\sum_j
\lambda_jh_j.
}
\]

Under suitable regularity, a local minimizer satisfies the KKT
conditions:

\[
\boxed{
\nabla_{\mathbf{x}}\mathcal{L}=0,
}
\]

\[
\boxed{
g_i(\mathbf{x})\leq0,
}
\]

\[
\boxed{
h_j(\mathbf{x})=0,
}
\]

\[
\boxed{
\mu_i\geq0,
}
\]

and

\[
\boxed{
\mu_i g_i(\mathbf{x})=0.
}
\]

Second-order conditions use

\[
\boxed{
\nabla_{\mathbf{x}\mathbf{x}}^2
\mathcal{L}
}
\]

restricted to feasible or critical directions.

Important numerical methods include:

\[
\boxed{
\text{penalty methods},
}
\]

\[
\boxed{
\text{augmented Lagrangians},
}
\]

\[
\boxed{
\text{SQP},
}
\]

and

\[
\boxed{
\text{interior-point methods}.
}
\]

Nonlinear programming therefore combines

\[
\boxed{
\text{calculus}
+
\text{constraint geometry}
+
\text{linear algebra}
+
\text{numerical methods}
+
\text{optimization theory}.
}
\]

%%%%%%%%%%%%%%%%%%%%%%%%%%%%%%%%%%%%%%%%%%%%%%%%%%%%%%%%%%%%
\subsection{Section Connections}
\label{subsec:nonlinear-programming-connections}
%%%%%%%%%%%%%%%%%%%%%%%%%%%%%%%%%%%%%%%%%%%%%%%%%%%%%%%%%%%%

\chapterconnection{
Nonlinear Programming extends the linear and general optimization
framework of Sections~\ref{sec:optimization-fundamentals} and
\ref{sec:linear-programming} to smooth and nonsmooth nonlinear models.
KKT systems, second-order conditions, penalty methods, augmented
Lagrangians, SQP, and nonlinear interior-point methods provide the central
tools for constrained nonlinear problems. The next section, Convex
Optimization, identifies an especially important subclass of nonlinear
optimization in which convexity converts local optimality into global
optimality and allows powerful duality and algorithmic guarantees.
}

%%%%%%%%%%%%%%%%%%%%%%%%%%%%%%%%%%%%%%%%%%%%%%%%%%%%%%%%%%%%
% SECTION SOURCE TRACKING
%%%%%%%%%%%%%%%%%%%%%%%%%%%%%%%%%%%%%%%%%%%%%%%%%%%%%%%%%%%%

%------------------------------------------------------------
% 10.3 Nonlinear Programming
%------------------------------------------------------------
%
% Source basis:
%   - Standard nonlinear programming theory.
%   - Standard constrained optimization theory.
%   - Standard numerical nonlinear optimization methods.
%
% Used for:
%   - nonlinear programming definitions;
%   - smooth and nonsmooth problems;
%   - local and global optimization;
%   - nonconvexity;
%   - feasible and tangent directions;
%   - active constraints;
%   - equality-constrained optimization;
%   - Lagrange multipliers;
%   - constraint qualifications;
%   - LICQ and MFCQ;
%   - nonlinear KKT conditions;
%   - complementary slackness;
%   - necessary and sufficient conditions;
%   - Hessian of the Lagrangian;
%   - reduced Hessians;
%   - Newton KKT systems;
%   - saddle-point systems;
%   - quadratic and exact penalty methods;
%   - barrier methods;
%   - central paths;
%   - augmented Lagrangians;
%   - SQP;
%   - quasi-Newton SQP;
%   - merit and filter methods;
%   - trust regions;
%   - feasibility restoration;
%   - nonlinear interior-point methods;
%   - active-set and projected methods;
%   - nonlinear least squares;
%   - Gauss--Newton;
%   - Levenberg--Marquardt;
%   - globalization;
%   - multistart and global optimization preview;
%   - scaling and sparse structure;
%   - automatic and finite-difference derivatives;
%   - KKT residuals;
%   - sensitivity analysis;
%   - engineering, statistical, economic, and scientific applications;
%   - exercises.
%
% Notes:
%   - Convex Optimization is developed next in Section 10.4.
%   - Integer Programming follows in Section 10.5.
%   - Numerical Optimization is developed more deeply in Chapter 11.
%
%%%%%%%%%%%%%%%%%%%%%%%%%%%%%%%%%%%%%%%%%%%%%%%%%%%%%%%%%%%%